
\documentclass{article}
\usepackage{colortbl}
\usepackage{makecell}
\usepackage{multirow}
\usepackage{supertabular}

\begin{document}

\newcounter{utterance}

\centering \large Interaction Transcript for game `wordle', experiment `high\_frequency\_words\_no\_clue\_no\_critic', episode 9 with deepseek{-}v4{-}pro.
\vspace{24pt}

{ \footnotesize  \setcounter{utterance}{1}
\setlength{\tabcolsep}{0pt}
\begin{supertabular}{c@{$\;$}|p{.15\linewidth}@{}p{.15\linewidth}p{.15\linewidth}p{.15\linewidth}p{.15\linewidth}p{.15\linewidth}}
    \# & \multicolumn{2}{c}{Player} && \multicolumn{2}{c}{Game Master} \\
    \hline

    \theutterance \stepcounter{utterance}  
    & & & \multicolumn{4}{p{0.6\linewidth}}{
        \cellcolor[rgb]{0.9,0.9,0.9}{
            \makecell[{{p{\linewidth}}}]{
                \texttt{\tiny{[P1$\langle$GM]}}
                \texttt{Jesteś czarodziejem języka, który lubi zgadywać słowa, korzystając z podanych zasad.} \\
\\ 
\texttt{Witamy w Wordle! Masz sześć prób, aby odgadnąć słowo docelowe, poprawne polskie słowo składające się z pięciu małych liter. Użyj tagów "wyjaśnienie:" oraz "zgadnij:", aby podać zwięzłe wyjaśnienie dla każdej próby.} \\
\\ 
\texttt{Na przykład, jeśli Twoją próbą jest "dzień", Twoja odpowiedź powinna brzmieć} \\
\texttt{wyjaśnienie: to jest popularne pięcioliterowe polskie słowo i zaczynam zgadywanie od tego słowa.} \\
\texttt{zgadnij: dzień} \\
\\ 
\texttt{Po każdej próbie Twoja odpowiedź zostanie zweryfikowana i otrzymasz informację zwrotną wskazującą, które litery są poprawne (zielony), które litery są poprawne, ale na niewłaściwej pozycji (żółty), oraz które litery są niepoprawne (czerwony). Ta informacja zwrotna może być przydatna przy ustalaniu, które litery uwzględnić lub wykluczyć w następnej próbie.} \\
\\ 
\texttt{Na przykład informacja zwrotna dla "dzień" może wyglądać tak:} \\
\texttt{informacja zwrotna: d<żółty> z<żółty> i<zielony> e<żółty> ń<czerwony>} \\
\\ 
\texttt{Wyjaśnienie powinno zawierać szczegóły dotyczące tego, jak wykorzystano informacja zwrotna, aby dojść do nowej próby.} \\
\\ 
\texttt{Nie generuj żadnego innego tekstu.} \\
\\ 
\texttt{Zacznijmy od Twojej pierwszej próby.} \\
            }
        }
    }
     \\ \\

    \theutterance \stepcounter{utterance}  
    & \multicolumn{4}{p{0.6\linewidth}}{
        \cellcolor[rgb]{0.561,0.737,0.561}{
            \makecell[{{p{\linewidth}}}]{
                \texttt{\tiny{[P1$\rangle$GM]}}
                \texttt{wyjaśnienie: Na początek wybieram popularne, neutralne słowo z różnymi samogłoskami i spółgłoskami, aby szybko sprawdzić występowanie liter.} \\
\texttt{zgadnij: serce} \\
            }
        }
    }
    & & \\ \\

    \theutterance \stepcounter{utterance}  
    & & & \multicolumn{2}{p{0.3\linewidth}}{
        \cellcolor[rgb]{0.9,0.9,0.9}{
            \makecell[{{p{\linewidth}}}]{
                \texttt{\tiny{[GM$|$GM]}}
                \texttt{attempts = 1} \\
\texttt{target\_word = klasa} \\
\texttt{guess = serce} \\
\texttt{guess\_feedback = s<yellow> e<red> r<red> c<red> e<red>} \\
            }
        }
    }
    & & \\ \\

    \theutterance \stepcounter{utterance}  
    & & & \multicolumn{4}{p{0.6\linewidth}}{
        \cellcolor[rgb]{0.9,0.9,0.9}{
            \makecell[{{p{\linewidth}}}]{
                \texttt{\tiny{[P1$\langle$GM]}}
                \texttt{informacja zwrotna: s<żółty> e<czerwony> r<czerwony> c<czerwony> e<czerwony>} \\
\\ 
\texttt{Podaj odpowiedź tylko w tym formacie.} \\
\texttt{wyjaśnienie: szczegóły} \\
\texttt{zgadnij: słowo} \\
            }
        }
    }
     \\ \\

    \theutterance \stepcounter{utterance}  
    & \multicolumn{4}{p{0.6\linewidth}}{
        \cellcolor[rgb]{0.561,0.737,0.561}{
            \makecell[{{p{\linewidth}}}]{
                \texttt{\tiny{[P1$\rangle$GM]}}
                \texttt{wyjaśnienie: Litera "s" jest w słowie, ale na innej pozycji. Litery "e", "r", "c" są wykluczone. Potrzebuję nowego słowa zawierającego "s" na innej pozycji, z nowymi samogłoskami, np. "a", "i" i niesprawdzonymi spółgłoskami.} \\
\texttt{zgadnij: sonda} \\
            }
        }
    }
    & & \\ \\

    \theutterance \stepcounter{utterance}  
    & & & \multicolumn{2}{p{0.3\linewidth}}{
        \cellcolor[rgb]{0.9,0.9,0.9}{
            \makecell[{{p{\linewidth}}}]{
                \texttt{\tiny{[GM$|$GM]}}
                \texttt{attempts = 2} \\
\texttt{target\_word = klasa} \\
\texttt{guess = sonda} \\
\texttt{guess\_feedback = s<yellow> o<red> n<red> d<red> a<green>} \\
            }
        }
    }
    & & \\ \\

    \theutterance \stepcounter{utterance}  
    & & & \multicolumn{4}{p{0.6\linewidth}}{
        \cellcolor[rgb]{0.9,0.9,0.9}{
            \makecell[{{p{\linewidth}}}]{
                \texttt{\tiny{[P1$\langle$GM]}}
                \texttt{informacja zwrotna: s<żółty> o<czerwony> n<czerwony> d<czerwony> a<zielony>} \\
\\ 
\texttt{Podaj odpowiedź tylko w tym formacie.} \\
\texttt{wyjaśnienie: szczegóły} \\
\texttt{zgadnij: słowo} \\
            }
        }
    }
     \\ \\

    \theutterance \stepcounter{utterance}  
    & \multicolumn{4}{p{0.6\linewidth}}{
        \cellcolor[rgb]{0.561,0.737,0.561}{
            \makecell[{{p{\linewidth}}}]{
                \texttt{\tiny{[P1$\rangle$GM]}}
                \texttt{wyjaśnienie: Litera "a" jest na właściwej, piątej pozycji. Litera "s" nadal jest żółta, więc nie jest na pierwszej ani na piątej pozycji. Wykluczone litery to e, r, c, o, n, d. Szukam słowa kończącego się na "a", zawierającego "s" na drugiej, trzeciej lub czwartej pozycji, np. "miska" lub "pasza", ale muszę unikać wykluczonych liter.} \\
\texttt{zgadnij: pasja} \\
            }
        }
    }
    & & \\ \\

    \theutterance \stepcounter{utterance}  
    & & & \multicolumn{2}{p{0.3\linewidth}}{
        \cellcolor[rgb]{0.9,0.9,0.9}{
            \makecell[{{p{\linewidth}}}]{
                \texttt{\tiny{[GM$|$GM]}}
                \texttt{attempts = 3} \\
\texttt{target\_word = klasa} \\
\texttt{guess = pasja} \\
\texttt{guess\_feedback = p<red> a<yellow> s<yellow> j<red> a<green>} \\
            }
        }
    }
    & & \\ \\

    \theutterance \stepcounter{utterance}  
    & & & \multicolumn{4}{p{0.6\linewidth}}{
        \cellcolor[rgb]{0.9,0.9,0.9}{
            \makecell[{{p{\linewidth}}}]{
                \texttt{\tiny{[P1$\langle$GM]}}
                \texttt{informacja zwrotna: p<czerwony> a<żółty> s<żółty> j<czerwony> a<zielony>} \\
\\ 
\texttt{Podaj odpowiedź tylko w tym formacie.} \\
\texttt{wyjaśnienie: szczegóły} \\
\texttt{zgadnij: słowo} \\
            }
        }
    }
     \\ \\

    \theutterance \stepcounter{utterance}  
    & \multicolumn{4}{p{0.6\linewidth}}{
        \cellcolor[rgb]{0.561,0.737,0.561}{
            \makecell[{{p{\linewidth}}}]{
                \texttt{\tiny{[P1$\rangle$GM]}}
                \texttt{wyjaśnienie: Piąta litera to "a". "s" jest żółte, więc nie jest na trzeciej pozycji. "a" jest żółte na drugiej pozycji, więc może być na pierwszej, trzeciej lub czwartej. Wykluczone litery: e, r, c, o, n, d, p, j. Brakujące litery to "s" oraz drugie "a". Popularne słowo to "klasa", ale "a" powtórzone byłoby na pozycjach 2 i 5, a "s" na 3? Sprawdzam: k l a s a {-}> a na 3 i 5, s na 4. To pasuje do dotychczasowych wskazówek.} \\
\texttt{zgadnij: klasa} \\
            }
        }
    }
    & & \\ \\

    \theutterance \stepcounter{utterance}  
    & & & \multicolumn{2}{p{0.3\linewidth}}{
        \cellcolor[rgb]{0.9,0.9,0.9}{
            \makecell[{{p{\linewidth}}}]{
                \texttt{\tiny{[GM$|$GM]}}
                \texttt{attempts = 4} \\
\texttt{target\_word = klasa} \\
\texttt{guess = klasa} \\
\texttt{guess\_feedback = k<green> l<green> a<green> s<green> a<green>} \\
            }
        }
    }
    & & \\ \\

    \theutterance \stepcounter{utterance}  
    & & & \multicolumn{2}{p{0.3\linewidth}}{
        \cellcolor[rgb]{0.9,0.9,0.9}{
            \makecell[{{p{\linewidth}}}]{
                \texttt{\tiny{[GM$|$GM]}}
                \texttt{game\_result = WIN} \\
            }
        }
    }
    & & \\ \\

\end{supertabular}
}

\end{document}
